\documentclass[12pt]{report}
\usepackage[utf8]{inputenc}
\usepackage[english, russian]{babel}
\usepackage{listings}
\usepackage{graphicx}
\usepackage{float}
\graphicspath{{imgs/}}
\usepackage{amsmath,amsfonts,amssymb,amsthm,mathtools} 
\usepackage{pgfplots}
\usepackage{filecontents}
\usepackage{indentfirst}
\usepackage{eucal}
\usepackage{enumitem}
\frenchspacing

\usepackage{indentfirst} % Красная строка

\usetikzlibrary{datavisualization}
\usetikzlibrary{datavisualization.formats.functions}

\usepackage{amsmath}
\usepackage{fixltx2e}
\usepackage{caption}


\definecolor{bluekeywords}{rgb}{0,0,1}
\definecolor{greencomments}{rgb}{0,0.5,0}
\definecolor{redstrings}{rgb}{0.64,0.08,0.08}
\definecolor{xmlcomments}{rgb}{0.5,0.5,0.5}
\definecolor{types}{rgb}{0.17,0.57,0.68}

\usepackage{listings}
\lstset{language=[Sharp]C,
	captionpos=t,
	numbers=left, %Nummerierung
	numberstyle=\small, % kleine Zeilennummern
	frame=single, % Oberhalb und unterhalb des Listings ist eine Linie
	stepnumber=1,                   
	numbersep=5pt,                
	showspaces=false,
	tabsize=2,
	showtabs=false,
	breaklines=true,
	showstringspaces=false,
	breakatwhitespace=true,
	escapeinside={(*@}{@*)},
	commentstyle=\color{greencomments},
	morekeywords={partial, var, value, get, set},
	keywordstyle=\color{bluekeywords},
	stringstyle=\color{redstrings},
	basicstyle=\ttfamily\small,
}

\usepackage[left=2cm,right=2cm, top=2cm,bottom=2cm,bindingoffset=0cm]{geometry}
% Для измененных титулов глав:
\usepackage{titlesec, blindtext, color} % подключаем нужные пакеты
\definecolor{gray75}{gray}{0.75} % определяем цвет
\newcommand{\hsp}{\hspace{20pt}} % длина линии в 20pt
% titleformat определяет стиль
\titleformat{\chapter}[hang]{\Huge\bfseries}{\thechapter\hsp\textcolor{gray75}{|}\hsp}{0pt}{\Huge\bfseries}

\usepackage{array}
\newcommand{\head}[2]{\multicolumn{1}{>{\centering\arraybackslash}p{#1}}{#2}}

% plot
\usepackage{pgfplots}
\usepackage{filecontents}
\usetikzlibrary{datavisualization}
\usetikzlibrary{datavisualization.formats.functions}

\begin{document}
%\def\chaptername{} % убирает "Глава"
\thispagestyle{empty}
\begin{titlepage}
	\noindent \begin{minipage}{0.15\textwidth}
		\includegraphics[width=\linewidth]{b_logo}
	\end{minipage}
	\noindent\begin{minipage}{0.9\textwidth}\centering
		\textbf{Министерство науки и высшего образования Российской Федерации}\\
		\textbf{Федеральное государственное бюджетное образовательное учреждение высшего образования}\\
		\textbf{~~~«Московский государственный технический университет имени Н.Э.~Баумана}\\
		\textbf{(национальный исследовательский университет)»}\\
		\textbf{(МГТУ им. Н.Э.~Баумана)}
	\end{minipage}
	
	\noindent\rule{18cm}{3pt}
	\newline\newline
	\noindent ФАКУЛЬТЕТ $\underline{\text{«Информатика и системы управления»}}$ \newline\newline
	\noindent КАФЕДРА $\underline{\text{«Программное обеспечение ЭВМ и информационные технологии»}}$\newline\newline\newline\newline\newline\newline\newline\newline\newline\newline\newline
	
	
	\begin{center}
		\noindent\begin{minipage}{1.3\textwidth}\centering
			\Large\textbf{  Отчет по лабораторной работе №7}\newline
			\textbf{по дисциплине \newline "Функциональное и логическое программирование"}\newline\newline
		\end{minipage}
	\end{center}
	
	\noindent\textbf{Тема} $\underline{\text{Рекурсивные функции}}$\newline\newline
	\noindent\textbf{Студент} $\underline{\text{Малышев И. А.}}$\newline\newline
	\noindent\textbf{Группа} $\underline{\text{ИУ7-61Б}}$\newline\newline
	\noindent\textbf{Оценка (баллы)} $\underline{\text{~~~~~~~~~~~~~~~~~~~~~~~~~~~}}$\newline\newline
	\noindent\textbf{Преподаватель: } $\underline{\text{Толпинская Н. Б.}}$\newline\newline\newline
	
	\begin{center}
		\vfill
		Москва~---~\the\year
		~г.
	\end{center}
\end{titlepage}


\setcounter{page}{2}

\chapter*{Практические задания}

\subsection*{1. Написать хвостовую рекурсивную функцию my-reverse, которая развернет верхний
	уровень своего списка-аргумента lst.}

\begin{lstlisting}[label=6xd, caption=Решение задания №1, language=lisp]
(defun my-reverse (lst)
	(my-rev lst ()))
	
(defun my-rev (lst acc)
	(cond ((null lst) acc)
		  (t (my-rev (cdr lst) (cons (car lst) acc)))))

\end{lstlisting}

\subsection*{2. Написать функцию, которая возвращает первый элемент списка-аргумента, который сам
	является непустым списком.}

\begin{lstlisting}[label=6xd, caption=Решение задания №2, language=lisp]
(defun find-first-lst (lst) 
	(cond ((null lst) nil)
		  ((listp (car lst)) (car lst)) 
		  (t (find-first-lst (cdr lst)))))

\end{lstlisting}

\subsection*{3. Написать функцию, которая выбирает из заданного списка только те числа, которые
	больше 1 и меньше 10.}

\begin{lstlisting}[label=6xd, caption=Решение задания №3, language=lisp]
(defun a1b10 (x) (and (> x 1) (< x 10)))

(defun a1b10-lst (lst acc) 
	(cond ((null lst) acc) 
		((a1b10 (car lst)) (a1b10-lst (cdr lst) (cons (car lst) acc))) 
		(t (a1b10-lst (cdr lst) acc))))
		
(defun above-one-below-ten (lst)
	(my-reverse (a1b10-lst lst acc)))

\end{lstlisting}

\subsection*{4. Напишите рекурсивную функцию, которая умножает на заданное число-аргумент все
	числа
	из заданного списка-аргумента, когда}

a) все элементы списка --- числа,

6) элементы списка -- любые объекты.

\begin{lstlisting}[label=6xd, caption=Решение задания №4, language=lisp]
; a)
(defun nums-mul (lst mul) 
	(cond ((null lst) nil)
		  (t (cons (* (car lst) mul) (nums-mul (cdr lst) mul)))))

; b)
(defun nums-mul (lst mul)
	(cond ((null lst) nil)
		  ((numberp (car lst)) (cons (* (car lst) mul) (nums-mul (cdr lst) mul)))
		  ((listp (car lst)) (cons (nums-mul (car lst) mul) (nums-mul (cdr lst) mul)))
		  (t (cons (car lst) (nums-mul (cdr lst) mul)))))

\end{lstlisting}

\subsection*{5. Напишите функцию, select-between, которая из списка-аргумента, содержащего только
	числа, выбирает только те, которые расположены между двумя указанными границами-аргументами и возвращает их в виде списка (упорядоченного по возрастанию списка чисел
	(+ 2 балла)).}

\begin{lstlisting}[label=6xd, caption=Решение задания №5, language=lisp]
(defun betweenp (el b1 b2)
	(< (* (- el b1) (- el b2)) 0))

(defun select-between (lst b1 b2)
	(cond ((null lst) nil)
			((betweenp (car lst) b1 b2) (cons (car lst) (select-between (cdr lst) b1 b2)))
			(T (select-between (cdr lst) b1 b2))))
	
\end{lstlisting}

\newpage
\subsection*{6. Написать рекурсивную версию (с именем rec-add) вычисления суммы чисел заданного
	списка:}

а) одноуровнего смешанного,

б) структурированного.

""\newline
\begin{lstlisting}[label=6xd, caption=Решение задания №6, language=lisp]
; a)
(defun rec-add-internal (lst sum) 
	(cond ((null lst) sum) 
		  (t (rec-add-internal (cdr lst) (+ sum (car lst))))))
		  
(defun rec-add (lst) 
	(rec-add-internal lst 0))


; b)
(defun rec-add-internal (lst sum) 
	(cond ((null lst) sum) 
		  ((numberp (car lst)) (rec-add-internal (cdr lst) (+ sum (car lst))))
		  ((listp (car lst)) (rec-add-internal (cdr lst) (rec-add-internal (car lst) sum)))
		  (t (rec-add-internal (cdr lst) sum))))

(defun rec-add (lst) 
	(rec-add-internal lst 0))
\end{lstlisting}

\subsection*{7. Написать рекурсивную версию с именем recnth функции nth.}

\begin{lstlisting}[label=6xd, caption=Решение задания №7, language=lisp]	
(defun recnth (lst n) 
	(cond ((null lst) nil) 
		  ((< n 0) nil) 
		  ((= n 0) (car lst)) 
		  (t (recnth (cdr lst) (- n 1)))))
	
\end{lstlisting}

\newpage
\subsection*{8. Написать рекурсивную функцию allodd, которая возвращает t когда все элементы списка
	нечетные.}

\begin{lstlisting}[label=6xd, caption=Решение задания №8, language=lisp]
(defun allodd (lst) 
	(cond ((null lst) nil) 
		  ((null (cdr lst)) (oddp (car lst))) 
		  ((oddp (car lst)) (allodd (cdr lst))) 
		  (t nil)))
	
\end{lstlisting}

\subsection*{9. Написать рекурсивную функцию, которая возвращает первое нечетное число из списка
	(структурированного), возможно создавая некоторые вспомогательные функции.}

\begin{lstlisting}[label=6xd, caption=Решение задания №9, language=lisp]
(defun first-odd (lst)
	(cond ((null lst) nil) 
			((numberp (car lst)) (cond ((oddp (car lst)) (car lst)) 
										(t (fodd (cdr lst)))))
			((listp (car lst)) (let ((tmp (fodd (car lst)))) 
									(cond ((null tmp) (fodd (cdr lst))) 
											(t tmp))))
			(T (fodd (cdr lst)))))

\end{lstlisting}

\subsection*{10. Используя cons-дополняемую рекурсию с одним тестом завершения,
	написать функцию которая получает как аргумент список чисел, а возвращает список
	квадратов этих чисел в том же порядке.}

\begin{lstlisting}[label=6xd, caption=Решение задания №10, language=lisp]
(defun square-list (lst)
	(cond ((null lst) lst)
		  ((atom lst) (* lst lst))
		  (T (cons (square-list (car lst))
				   (square-list (cdr lst))))))
	
\end{lstlisting}


\end{document}